\documentclass[11pt,a4paper]{article}
\usepackage[T1]{fontenc}
\usepackage{lmodern}
\usepackage[margin=26mm]{geometry}
\usepackage{amsmath,amssymb,amsthm,booktabs,array,longtable}
\usepackage{microtype}
\usepackage[hidelinks]{hyperref}
\hypersetup{pdftitle={Further OEIS proofs and formalizations},pdfauthor={Dev Goyal},pdfsubject={AI-generated proofs, Lean 4 formalization, and precise prior-art accounting}}
\setlength{\parskip}{3pt}
\setlength{\emergencystretch}{2em}
\newtheorem{theorem}{Theorem}[section]
\newtheorem{lemma}[theorem]{Lemma}
\newtheorem{corollary}[theorem]{Corollary}
\theoremstyle{remark}
\newtheorem{remark}[theorem]{Remark}
\newcommand{\Z}{\mathbb Z}
\newcommand{\Q}{\mathbb Q}
\newcommand{\N}{\mathbb N}
\newcommand{\F}{\mathbb F}
\newcommand{\ps}[1]{#1[[x]]}
\newcommand{\ord}{\operatorname{ord}_x}
\title{Further OEIS proofs and formalizations\\\large Expanded proof package, version 2}
\author{Dev Goyal}
\date{24 October 2026}
\begin{document}
\maketitle

\begin{center}
\fbox{\begin{minipage}{0.90\linewidth}
\textbf{AI disclosure and review status.}
ChatGPT (OpenAI) generated the new proofs, exposition, Lean formalization,
verification programs, and AI-assisted reviews in response to Dev Goyal's
requests. Dev Goyal is the sole human author and initiated the work.
\textbf{Independent human review is pending.} Upstream code and earlier
mathematical results are credited separately. This document does not claim
OEIS acceptance, peer review, or exclusive priority.
\end{minipage}}
\end{center}

\begin{abstract}
We give formally checked proofs of seven statements labeled as conjectures
in four retrieved OEIS entries: two parity assertions for A396093, two
congruences each for A389474 and A389475, and the coefficient congruence
for A396102. Two of the cubic assertions are direct instances of the
standard Frobenius identity. The sequence-specific arguments give a complete
period-ten parity classification, an identity for an entire square modulo
nine, and a general family of inverse-iterate congruences. We construct the
implicit integer series and prove their source equations and uniqueness.
We also formalize a totient-period theorem for all factorial-weighted
Stirling transforms and seven named applications. The latter are already
covered by Danesh's 2026 theorem, with a stronger Carmichael bound, and
are not counted as newly solved open problems. The archive retains all
previous proofs, the previous manuscript, and the previous ZIP verbatim.
\end{abstract}

\section{Scope, notation, and historical status}

All series in this paper are formal. A normalized series means
$U=x+O(x^2)$. We distinguish ordinary powers $U^k$ from compositional
iterates $U^{\circ k}$, with $U^{\circ0}=x$. A congruence of series is
coefficientwise, and $\N$ includes zero. In particular, no convergence
assumption is required for any statement below.

The source pages were consulted on 24 October 2026. The four entries in
Table~\ref{tab:new} explicitly labeled the indicated statements as
conjectures \cite{oeis093,oeis474,oeis475,oeis102}. That label is evidence
about the retrieved entry, not a guarantee that no proof exists elsewhere.
Our contribution is a reproducible proof package with exact mathematical
scope. We make no global novelty claim. The two Frobenius consequences
are standard algebra applied to the stated sequences.

\begin{table}[ht]
\centering\small
\begin{tabular}{@{}lp{0.68\linewidth}r@{}}
\toprule
Entry & Statements addressed & Count\\
\midrule
A396093 & Even coefficients at even indices; exact condition for evenness at odd indices & 2\\
A389474 & Square coefficients above degree three vanish modulo nine; cube-power congruence modulo three & 2\\
A389475 & Cube substitution and the associated square identity modulo three & 2\\
A396102 & Every positive-index coefficient is one modulo three & 1\\
\bottomrule
\end{tabular}
\caption{Seven newly added proofs of statements marked conjectural in the retrieved entries.}\label{tab:new}
\end{table}

The seven additional periodicity applications in Section~\ref{sec:stirling}
have a different historical status. Danesh's Theorem~4.1 and
Corollary~4.2 already imply them \cite{danesh}. They are included as
formalizations and recovered corollaries of known mathematics. The package
does not inflate its open-problem count by treating those seven cases,
their general theorem, or its intermediate lemmas as new discoveries.

\section{A396093: a complete parity classification}

The original definition is
\[
B(x)=\frac{x}{(1-x)^2},\qquad A=B\circ B\circ B,
\qquad a_n=[x^n]A\in\Z.
\]
The offset is zero, so $a_0=0$. These definitions, rather than a conjectured
coefficient pattern, are the starting point of the Lean construction.

\begin{theorem}\label{thm:parity}
For every $n\ge0$,
\[
a_n\equiv
\begin{cases}
1\pmod2,&n\equiv1,3,7,9\pmod{10},\\
0\pmod2,&\text{otherwise}.
\end{cases}
\]
Consequently $a_{2n}$ is even for every $n\ge0$. For $n\ge1$,
\[
a_{2n-1}\text{ is even}
\quad\Longleftrightarrow\quad n\equiv3\pmod5
\quad\Longleftrightarrow\quad n=5k-2\text{ for some }k\ge1.
\]
\end{theorem}

\begin{proof}
Reduce coefficients to $\F_2$. In characteristic two,
$B(x)=x/(1+x^2)$. If $T=P/D$ has zero constant coefficient and $D(0)=1$,
then
\[
B(T)=\frac{PD}{D^2+P^2}.
\]
Every denominator has constant coefficient one and is therefore an
invertible formal series. Applying this identity gives
\[
\begin{array}{c|c|c}
j & \text{numerator of }B^{\circ j} & \text{denominator}\\ \hline
1 & x & 1+x^2\\
2 & x+x^3 & 1+x^2+x^4\\
3 & x+x^7 & 1+x^2+x^4+x^6+x^8.
\end{array}
\]
Multiplication of the last numerator and denominator by $1+x^2$ gives
the exact identity
\begin{equation}\label{eq:parityrational}
A(x)(1+x^{10})=x+x^3+x^7+x^9\qquad\text{in }\ps{\F_2}.
\end{equation}
Since $(1+x^{10})^{-1}=\sum_{j\ge0}x^{10j}$ in this ring,
\eqref{eq:parityrational} yields precisely the four stated residue classes.
Substitute $2n$ and $2n-1$ to obtain the two original assertions. For
the latter, $2n-1$ is the omitted odd residue class $5$ modulo ten
exactly when $n\equiv3\pmod5$.
\end{proof}

The formal proof works with the integral series $x(\sum_{n\ge0}x^n)^2$
and its actual third composition. It proves that reduction commutes with
composition. The denominator cancellation in
\eqref{eq:parityrational} uses a constructed unit inverse; it does not
assume cancellation in an arbitrary ring. The final Lean statements
retain the source's $n\ge1$ and $k\ge1$ restrictions where subtraction
of natural numbers occurs.

\section{The two cubic series: existence before congruences}

\subsection{A canonical integral family}

\begin{theorem}\label{thm:cubicexist}
For every $s\in\Z$, there is a unique normalized series
$R_s\in\ps{\Z}$ satisfying
\begin{equation}\label{eq:cubicfamily}
R_s(3x^2)=2x^2+R_s(x)^2+2sR_s(x)^3.
\end{equation}
\end{theorem}

\begin{proof}
For any normalized integral series $U$, define its residual by
\[
E_s(U)=U^2+2sU^3-U(3x^2)+2x^2.
\]
Every coefficient of $E_s(U)$ is even. Indeed, after reduction modulo
two, Frobenius gives $U(x)^2=U(x^2)$ and $3x^2=x^2$, so the residual
vanishes. Also $E_s(U)=O(x^3)$ for every normalized $U$.

Suppose $V-U=cx^r+O(x^{r+1})$, with $r\ge2$. The quadratic difference
has leading term $2cx^{r+1}$; the cubic difference has order at least
$r+2$; and $(V-U)(3x^2)$ has order at least $2r\ge r+2$.
Consequently
\begin{equation}\label{eq:triangular}
E_s(V)-E_s(U)=2cx^{r+1}+O(x^{r+2}).
\end{equation}
Starting with $U=x$, choose its degree-$r$ coefficient to cancel the
degree-$(r+1)$ coefficient of the residual. The needed division by two
is exact by the preceding parity argument. Equation~\eqref{eq:triangular}
shows that earlier cancellations remain valid. Each degree is thus
determined by finitely many earlier coefficients, and the resulting
integral series has zero residual in every degree.

For uniqueness, take the first degree $r\ge2$ where two candidate
solutions differ. Equation~\eqref{eq:triangular} would give a nonzero
degree-$(r+1)$ difference of residuals, because $2c\ne0$ in $\Z$.
This contradicts both residuals being zero.
\end{proof}

The proof formalizes divisibility of the residual by two before using
integer division. Thus the construction does not assume the existence
or integrality of an OEIS sequence. The names used below denote the
uniquely constructed series.

\subsection{A389474 and a square-zero substitution}

Let $R=R_1$ and let $U$ be its compositional inverse. Such an inverse
exists over $\Z$ for every normalized series: its degree-$n$ coefficient
is determined with leading multiplier one, recursively. Both inverse
identities are proved in the reused formal-series infrastructure.
Substituting $U(x)$ in \eqref{eq:cubicfamily} gives
\[
R(3U(x)^2)=2U(x)^2+x^2+2x^3.
\]
Applying $U$ proves the original A389474 equation
\begin{equation}\label{eq:474}
U\bigl(x^2+2x^3+2U(x)^2\bigr)=3U(x)^2.
\end{equation}
Conversely, if a normalized $U$ satisfies \eqref{eq:474}, substitute its
inverse $R$ and then apply $R$ to recover \eqref{eq:cubicfamily} for
$s=1$. Therefore Theorem~\ref{thm:cubicexist} also gives existence and
uniqueness for the literal A389474 source equation.

\begin{theorem}\label{thm:modnine}
The A389474 series satisfies
\[
U(x)^2\equiv x^2+2x^3\pmod9.
\]
In particular, $[x^n]U(x)^2$ is divisible by nine for every $n>3$.
\end{theorem}

\begin{proof}
Work in $(\Z/9\Z)[[x]]$ and put $S=U^2$, $Q=x^2+2x^3$.
The source equation becomes $U(Q+2S)=3S$. Let $V$ be the genuine
normalized compositional inverse of $U$ over this ring. Apply $V$ to
obtain
\[
Q+2S=V(3S).
\]
The series $3S$ has zero constant coefficient and is square-zero:
$(3S)^2=9S^2=0$. A normalized series has the form
$V(t)=t+t^2H(t)$, so $V(3S)=3S$. It follows that $Q+2S=3S$, hence
$S=Q$, as required. No division by three is used.
\end{proof}

This establishes a full formal-series identity, including the low-degree
terms, and hence strengthens the source's statement restricted to
degrees above three.

\subsection{Frobenius and A389475}

For any $W\in\ps{\Z}$, prime $p$, and coefficientwise reduction
$\overline W\in\ps{\F_p}$, the standard Frobenius identity is
\begin{equation}\label{eq:frobenius}
\overline W(x)^p=\overline W(x^p).
\end{equation}
One may prove it degree by degree: truncate $W$ above the degree in
question, apply the polynomial identity in characteristic $p$, and
use $a^p=a$ in $\F_p$. The formal proof follows this finite-truncation
argument. In particular, for every $k\ge0$,
\[
U(x)^{3k}\equiv U(x^3)^k\pmod3.
\]
This proves A389474's second conjecture, including the otherwise
unneeded cases $k=0$ and coefficient index zero. It is a universal
algebraic fact, not a special property newly discovered for A389474.

Now let $V=R_{-1}$, which is the A389475 series. It satisfies
\begin{equation}\label{eq:475}
V(3x^2)=2x^2+V^2-2V^3.
\end{equation}
Frobenius directly proves the first requested identity
$V^3\equiv V(x^3)\pmod3$. Since $V(0)=0$, reduction of
\eqref{eq:475} modulo three gives
\[
0=2x^2+V^2-2V^3,
\]
or $V^2=x^2+2V^3$ in $\ps{\F_3}$. Substitution of the Frobenius
identity proves the second:
\begin{equation}\label{eq:475square}
V(x)^2\equiv x^2+2V(x^3)\pmod3.
\end{equation}
These are exactly the two assertions labeled (4.a) and (4.b) in the
retrieved A389475 entry. The formal development also supplies the
literal source-equation existence and uniqueness theorem.

\section{A396102 and the inverse-iterate family}

For each positive integer $r$, consider the normalized series equation
\begin{equation}\label{eq:inverseoriginal}
F_r^{\circ r}(x)=(1+x)F_r^{\circ(r-1)}(x).
\end{equation}
The case $r=3$ is the original A396102 equation. We construct it through
an auxiliary series, rather than assuming that a solution exists.

\begin{lemma}\label{lem:inverseconstruct}
For every $r\ge0$ and every commutative coefficient ring, there is a
unique normalized series $R_r$ satisfying
\begin{equation}\label{eq:inverseaux}
R_r=x-R_r R_r^{\circ r}.
\end{equation}
If $r\ge1$, its compositional inverse $F_r$ is the unique normalized
solution of \eqref{eq:inverseoriginal}.
\end{lemma}

\begin{proof}
Define $T_r(G)=x-GG^{\circ r}$. It preserves normalization. If two
normalized series agree through degree $n-1$, their iterates do also.
The difference of the two products $GG^{\circ r}$ and $HH^{\circ r}$
has one factor divisible by $x^n$ and another divisible by $x$ in
each term. Thus
\[
G-H=O(x^n)\quad\Longrightarrow\quad
T_r(G)-T_r(H)=O(x^{n+1}).
\]
Starting at $x$ and iterating $T_r$, every coefficient stabilizes after
finitely many steps. The stabilized coefficients give a fixed point;
the same order-increase implication proves uniqueness. This is a
purely coefficientwise construction and works over arbitrary
commutative rings.

Take the genuine compositional inverse $F_r$ of $R_r$. Substitute
$F_r^{\circ r}(x)$ into \eqref{eq:inverseaux}. Inverse cancellation
gives
\[
F_r^{\circ(r-1)}=F_r^{\circ r}-xF_r^{\circ(r-1)},
\]
which is \eqref{eq:inverseoriginal}. Conversely, substitute
$R^{\circ r}$ into \eqref{eq:inverseoriginal}, where $R=F^{-1}$.
This gives $x=(1+R^{\circ r})R$, hence \eqref{eq:inverseaux}.
Uniqueness of $R_r$ now implies uniqueness of $F_r$.
\end{proof}

\begin{theorem}\label{thm:inversemod}
For every $r\ge1$,
\[
R_r(x)\equiv\frac{x}{1+x}\pmod r,\qquad
F_r(x)\equiv\frac{x}{1-x}\pmod r.
\]
Consequently, for $n\ge1$ and $k\ge0$,
\begin{equation}\label{eq:iteratemod}
[x^n]F_r^{\circ k}(x)\equiv k^{n-1}\pmod r.
\end{equation}
Here $0^0=1$, consistently with the degree-one coefficient of the
zeroth compositional iterate $x$.
\end{theorem}

\begin{proof}
Write $M_c(x)=x/(1-cx)$. Direct substitution gives
$M_a\circ M_b=M_{a+b}$ over any commutative ring. Therefore
$M_{-1}^{\circ r}=M_{-r}=x$ over $\Z/r\Z$. Moreover
\[
x-M_{-1}(x)x=M_{-1}(x),
\]
so $M_{-1}$ satisfies the auxiliary fixed-point equation modulo $r$.
The coefficientwise uniqueness in Lemma~\ref{lem:inverseconstruct}
identifies it with the reduction of $R_r$. Its actual inverse is
$M_1$, identifying the reduction of $F_r$. Finally,
$M_1^{\circ k}=M_k=x/(1-kx)$, whose degree-$n$ coefficient is
$k^{n-1}$ for $n\ge1$.
\end{proof}

For $r=3$ and $k=1$, this proves the A396102 conjecture
\[
[x^n]F_3(x)\equiv1\pmod3\qquad(n\ge1).
\]
The formal proof includes both directions of the correspondence with
the original triple-composition equation, not merely a congruence for
an auxiliary series. The formal parameterization uses $r+1$ when
stating the original equation so that its predecessor iterate is
always a natural number. The coefficient theorem itself also has a
well-defined $r=0$ version in \texttt{ZMod 0}; no such extension is
needed to settle the stated OEIS problem.

\section{Known Stirling periodicity, now formalized}\label{sec:stirling}

Let $S(n,k)$ be the Stirling numbers of the second kind, defined by
$S(0,0)=1$, zero on the other boundary entries, and
\[
S(n+1,k+1)=(k+1)S(n,k+1)+S(n,k).
\]
For an arbitrary sequence of integers $c_0,c_1,\ldots$, put
\begin{equation}\label{eq:transform}
b_n=\sum_{k=0}^n c_k k!S(n,k).
\end{equation}
No positivity, boundedness, or recurrence condition is imposed on the
weights. The usual formal exponential identity in $\Q[[x]]$ identifies
\eqref{eq:transform} as the coefficient sequence of
\[
\sum_{n\ge0}b_n\frac{x^n}{n!}=G(e^x-1),\qquad
G(t)=\sum_{k\ge0}c_kt^k\in\Z[[t]].
\]

\begin{theorem}\label{thm:stirling}
For every $m>0$ and $n\ge m$,
\[
b_{n+\varphi(m)}\equiv b_n\pmod m.
\]
The least positive period of the tail beginning at index $m$ exists
and divides $\varphi(m)$.
\end{theorem}

\begin{proof}
We first show that every natural base $a$ satisfies
\begin{equation}\label{eq:powerperiod}
a^{n+\varphi(m)}\equiv a^n\pmod m\qquad(n\ge m).
\end{equation}
It suffices to check each prime power $p^e$ dividing $m$. If $p$ divides
$a$, then $e\le p^e\le m\le n$, so both powers vanish modulo $p^e$.
If $p$ does not divide $a$, Euler's theorem applies modulo $p^e$,
and $\varphi(p^e)$ divides $\varphi(m)$. This proves
\eqref{eq:powerperiod} by prime-power divisibility.

The power expansion is
\begin{equation}\label{eq:powerexp}
a^n=\sum_{k=0}^{a}\binom ak k!S(n,k)\qquad(a\le n).
\end{equation}
It follows from the Stirling recurrence and the usual recurrence for
falling factorials. Fix $n\ge m$. Strong induction on $k<m$ shows
that $k!S(n+\varphi(m),k)$ and $k!S(n,k)$ are congruent modulo $m$:
apply \eqref{eq:powerexp} at $a=k$ to the two exponents, cancel all
terms with index below $k$ using the induction hypothesis, and note
that $\binom kk=1$. This argument makes no division by $k!$ modulo
$m$.

For $k\ge m$, the factor $k!$ is divisible by $m$. Thus
\eqref{eq:transform} reduces to the fixed finite sum with $k<m$.
The preceding columnwise congruence proves the claimed period for
arbitrary integral weights. Since $\varphi(m)>0$, the tail has at
least one positive period.

For completeness, let $d$ be its least positive period and write any
period $p=qd+r$, $0\le r<d$. Repeated use of the $d$-periodicity
shows that $r$ is also a period. Minimality forces $r=0$, so $d$ divides
$p$. Apply this with $p=\varphi(m)$. In Lean this last step uses the
standard minimal-period theorem for the shift map on the entire tail.
\end{proof}

\paragraph{Prior work and exact formal scope.}
Danesh's 26 August 2026 preprint proves the general integral-substitution
result with the stronger eventual period $\lambda(m)$, the Carmichael
function \cite{danesh}. It therefore already covers Theorem~\ref{thm:stirling}
as an eventual-period assertion and all seven applications below.
Our displayed bound $n\ge m$ is explicit; no priority claim is made
for this elementary threshold. Schreyer earlier treated Fubini and
$r$-Fubini numbers with a Carmichael bound \cite{schreyer}.

The Lean theorem is the fully quantified finite-sum statement
\eqref{eq:transform}, with exact definitions for every named weight.
The identity involving the formal exponential series is explained
mathematically here and is \emph{not separately claimed as a Lean
theorem}. For the named applications, the OEIS entries provide the
same Stirling-transform formulas or equivalent partition generating
functions. Numerical prefix checks further corroborate their indexing.

\begin{table}[ht]
\centering\small
\begin{tabular}{@{}ll@{}}
\toprule
Entry & Weight $c_k$ in \eqref{eq:transform}\\
\midrule
A052895 & Catalan number $\frac1{k+1}\binom{2k}{k}$\\
A064618 & $k!$\\
A354242 & $\binom{2k}{k}$\\
A167137 & Number of integer partitions of $k$\\
A305550 & Number of partitions of $k$ into distinct parts\\
A306082 & Number of partitions of $k$ into square parts\\
A080253 & $1,3,7,17,\ldots$, with $c_{k+2}=2c_{k+1}+c_k$\\
\bottomrule
\end{tabular}
\caption{Seven credited applications of an already proved general theorem.}
\end{table}

The partition weights in Lean are actual finite cardinalities of
\texttt{Nat.Partition}, with distinctness or square-part restrictions
where appropriate. The Pell weight is an explicit recurrence. No named
weight is left as an unspecified function assumed to match an OEIS
sequence. We do not address the separate continued-fraction conjecture
in A354242 or the separate power/logarithm integrality conjecture in
A167137.

\section{Formal verification and reproducibility}

The combined package pins Lean 4.19.0 and Mathlib revision
\begin{center}\small\ttfamily
c44e0c8ee63ca166450922a373c7409c5d26b00b.
\end{center}
The additional results reside in \texttt{lean/More/}. The principal
module groups are:
\begin{itemize}
\item \texttt{RationalParity.lean}: the actual third composition,
characteristic-two denominator identity, and exact parity conclusions.
\item \texttt{CubicFrobenius.lean}, \texttt{CubicReversion.lean}, and
\texttt{CubicConsequences.lean}: Frobenius, integral construction,
source-equation equivalence, and all four cubic congruences.
\item \texttt{InverseIterate.lean}: the auxiliary fixed point,
genuine inverse, original source equation, uniqueness, and congruences.
\item The Stirling, power-periodicity, and tail-periodicity modules:
the finite transform theorem, actual named weights, and least periods.
\end{itemize}

The release checker has a fixed list of \textbf{130 theorem endpoints}:
the original 67 and 63 added endpoints, including construction and
transport results as well as named consequences. This is a proof-audit
count, not a count of open problems. The checker removes the first-party
build cache, rebuilds the project, compiles explicit semantic contracts,
prints fully elaborated theorem types, and checks every endpoint's
transitive axiom dependencies. The only permitted axioms are
\texttt{propext}, \texttt{Classical.choice}, and \texttt{Quot.sound}.
No \texttt{sorry}, \texttt{admit}, \texttt{native\_decide}, or custom
mathematical axiom is used in the proof sources.

\texttt{More/Contracts.lean} states the source equations, definitions,
coefficient domains, and unbounded conclusions independently of the
exported theorem types. The old contracts are retained unchanged within
the expanded contract file. Mechanical success establishes the stated
Lean propositions relative to Lean's foundations; checking the
translation from OEIS prose to those propositions is a separate
semantic responsibility. AI reviewers examined those correspondences,
and the package records their findings. Independent human review
remains pending.

\paragraph{Finite calculations.}
Independent exact-integer scripts construct coefficients from the
original equations, rather than from the desired modular patterns.
The cubic check runs through degree 40 and tests both equations through
degree 41. The inverse-iterate check covers parameters $1$ through $8$
and iterates $0$ through $8$ through degree 40. The parity check uses
direct composition, a separate coefficient formula, and the rational
recurrence, including parity through degree 2000. The Stirling check
covers moduli $1$ through $128$, signed as well as positive weights,
and 81,376 period comparisons. JSON files record their exact scopes.
These calculations are corroboration; the Lean proofs establish the
claims for every allowed index.

\paragraph{Reproduction.}
With the pinned Lean toolchain and Python 3.11 or later installed, run
the following from the extracted directory:
\begin{verbatim}
cd lean
lake exe cache get
lake build
cd ..
python3 scripts/check_formalization.py
\end{verbatim}
The checker rejects missing endpoints, changed dependency pins,
warning/error output, extra axioms, and first-party source changes
during validation. It also verifies the previous archive's fixed
SHA-256 and the unchanged previous Lean proof modules. The generated
report and logs in \texttt{reproducibility/lean\_complete/} are the
authoritative record of the delivered run.

\section{What the combined archive retains}

The previous package's compositional-congruence manuscript and all
of its Lean modules are included. Those cover all four A396099 claims;
the A396798 iterate assertions and the correction of a false starting
phase; A396807 including integer iterates; the A393856--A393859
scaled-inverse assertions; and the original A396843-family acceptance
theorems. Their complete proofs remain in the previous manuscript and
source. This new document adds to that work rather than revising its
mathematical conclusions.

For exact preservation, the archive also contains the original ZIP at
\begin{center}\small
\path{prior_work/oeis_congruence_proofs_Dev_Goyal_v1.zip}.
\end{center}
Its SHA-256 is
\begin{center}\small\ttfamily
04f347b0e116a1c92cd7a1f21f38a37ae3d\\
fb71c733dd5f55769accd0714489c
\end{center}
The line break is typographical; the checksum is one 64-character string
in the validation manifest. The package preserves the author details
from the supplied repository and does not reuse the earlier work's DOI
as a DOI for this expanded release.

The general Stirling recurrence and part of its power-expansion proof
were adapted from a later Apache-2.0 Mathlib module. Its copyright,
original authors, exact revision, license, and change notice are
retained under \texttt{third\_party/mathlib/}. This upstream reuse is
separate from the AI-generated modular transfer and sequence-specific
proofs.

\clearpage
\begin{thebibliography}{20}
\raggedright
\bibitem{oeis093} P. D. Hanna, OEIS A396093,
\url{https://oeis.org/A396093}. Accessed 24 October 2026.
\bibitem{oeis474} P. D. Hanna, OEIS A389474,
\url{https://oeis.org/A389474}. Accessed 24 October 2026.
\bibitem{oeis475} P. D. Hanna, OEIS A389475,
\url{https://oeis.org/A389475}. Accessed 24 October 2026.
\bibitem{oeis102} P. D. Hanna, OEIS A396102,
\url{https://oeis.org/A396102}. Accessed 24 October 2026.
\bibitem{danesh} P. Danesh, \emph{Eventual periodicity of weighted
Stirling transforms and two conjectures of Bala}, version 1,
26 August 2026, Theorem 4.1 and Corollary 4.2.
\url{https://doi.org/10.33774/coe-2026-4d6jr}.
Preprint; not peer-reviewed by Cambridge University Press at posting.
\bibitem{schreyer} B. Schreyer, \emph{Rigged Horse Numbers and their
Modular Periodicity}, arXiv:2409.03799v2, 3 December 2024.
\url{https://arxiv.org/abs/2409.03799}.
\bibitem{oeisstirling} OEIS entries A052895, A064618, A354242,
A167137, A305550, A306082, and A080253. Each is available at
\texttt{https://oeis.org/} followed by its identifier. Accessed
24 October 2026; precise formulas and status are in the source ledger.
\bibitem{repo} D. Goyal, \emph{oeis-a396843-proof}, supplied source
repository, revision \texttt{995e55adcbe06b2a913e24f12ff73a7faf362f06}.
\url{https://github.com/dvgyl/oeis-a396843-proof}.
\bibitem{mathlib} Mathlib contributors, formal proof library for Lean.
Pinned dependency and adapted upstream module revisions are recorded
in \texttt{lean/lake-manifest.json} and
\texttt{third\_party/mathlib/NOTICE.md}, respectively.
\end{thebibliography}
\end{document}
